\begin{center}
    \textbf{\large VỊ TRÍ ỨNG TUYỂN: CHUYÊN VIÊN KỸ THUẬT CÔNG NGHỆ SỐ VÀ CHUYÊN VIÊN TRUYỀN THÔNG, MARKETING}
\end{center}
\vspace{0.2cm}

\section{Giới thiệu}
\begin{onecolentry}
    \begin{highlights}
        \item Chuyên viên Tối ưu hóa Phòng vô tuyến và Kỹ sư Hệ thống với kinh nghiệm chuyên sâu về hạ tầng viễn thông, ảo hóa IDC và phát triển phần mềm hiệu năng cao.
        \item Năng lực Công nghệ số: Tự động hóa quy trình (CI/CD, DevOps), triển khai các hệ thống phân tích dữ liệu lớn và thiết kế kiến trúc hạ tầng mạng lưới, máy chủ (Server/Storage).
        \item Kinh nghiệm Tích hợp Truyền thông \& Kinh doanh: Xây dựng nền tảng xử lý đơn hàng đa kênh (Shopee, TikTok, Lazada), phát triển đường ống dữ liệu hướng sự kiện nhằm phân tích và nâng cao trải nghiệm khách hàng.
    \end{highlights}
\end{onecolentry}

\section{Học vấn}
\begin{twocolentry}{\textit{2019 -- 2023}}
    \textbf{Trường Đại học Bách Khoa TP.HCM - Đại học Quốc gia TP.HCM (VNU-HCM)}\\[4pt]
    \textit{Chuyên ngành Khoa học Máy tính}
\end{twocolentry}

\section{Ngoại ngữ}
\begin{onecolentry}
    \begin{highlights}
        \item Tiếng Anh, IELTS: 6.0
    \end{highlights}
\end{onecolentry}

\section{Kinh nghiệm Làm việc}

% --- VNPT (Hiện tại) ---
\begin{twocolentry}{\textit{Tháng 4 2026 -- Hiện tại}}
    \textbf{Công ty Công nghệ Thông tin VNPT (VNPT-IT)} \\[4pt]
    \textit{Kỹ sư DevSecOps \& Platform (Kubernetes / MLOps)}
\end{twocolentry}
\begin{onecolentry}
    \begin{highlights}
        \item Chịu trách nhiệm chính nền tảng \textbf{Kubernetes (RKE2)} production on-premise --- tổng \textbf{140 vCPU / hơn 400 GB RAM / NVIDIA A100 80GB}, vận hành \textbf{trên 100 pod} thuộc hơn 20 namespace, tỉ lệ workload khỏe mạnh 99\%.
        \item Xây dựng và vận hành \textbf{hơn 30 pipeline CI/CD} (Jenkins multibranch, GitLab CI) cho dịch vụ .NET, Go, Node.js và Python/AI; chuẩn hóa khuôn pipeline và tự động đồng bộ sang 18 repository qua GitLab API.
        \item \textbf{DevSecOps:} tích hợp quét bảo mật vào mọi pipeline --- \textbf{Trivy} (ảnh container), \textbf{gitleaks} (lộ bí mật), \textbf{SonarQube} (chất lượng mã); gia cố workload chạy không quyền root với \texttt{seccomp} và phân tách mạng bằng NetworkPolicy.
        \item \textbf{MLOps:} vận hành hạ tầng GPU (A100 80GB) và CI/CD cho các dịch vụ mô hình thị giác máy tính, cùng \textbf{MLflow} và luồng dữ liệu \textbf{Airflow / Spark / NiFi}.
        \item Triển khai \textbf{Debezium / Kafka Connect} bắt biến đổi dữ liệu (CDC) trên PostgreSQL logical replication; vận hành Kafka, Redis, Elasticsearch, MongoDB, PostGIS và MinIO.
        \item Truy vết nguyên nhân gốc sự cố production kéo dài nhiều giờ, sửa tận gốc trong mã nguồn, sau đó bổ sung \textbf{startup/readiness probe} và định mức tài nguyên để biến restart gián đoạn thành cập nhật cuốn chiếu không downtime.
    \end{highlights}
\end{onecolentry}
\vspace{0.2 cm}

% --- MobiFone ---
\begin{twocolentry}{\textit{Tháng 4 2025 -- Tháng 4 2026}}
    \textbf{Tổng công ty Viễn thông MobiFone} \\[4pt]
    \textit{Chuyên viên Tối ưu hóa Phòng vô tuyến và Kỹ sư Công nghệ thông tin}
\end{twocolentry}
\begin{onecolentry}
    \begin{highlights}
        \item Thiết kế, triển khai và tối ưu hóa hạ tầng CNTT doanh nghiệp tính sẵn sàng cao trên các thiết bị vật lý (máy chủ Dell, HPE) và các giải pháp lưu trữ toàn diện (SAN/NAS hỗ trợ phân tầng lưu trữ).
        \item Kiến trúc và quản trị các cụm ảo hóa quy mô lớn sử dụng VMware vSphere và Proxmox VE, triển khai các tính năng nâng cao như High Availability (HA), vMotion và Distributed Resource Scheduler (DRS) nhằm đảm bảo uptime 99.99\%.
        \item Dẫn dắt toàn bộ vòng đời tích hợp cho các dự án hạ tầng phức tạp: tiến hành khảo sát thực địa, soạn thảo tài liệu thiết kế kỹ thuật (HLD/LLD), và trực tiếp chỉ đạo các giai đoạn chạy thử nghiệm (pilot) cùng nghiệm thu (UAT) trước khi bàn giao.
        \item Quản trị an toàn hệ thống Microsoft Enterprise tập trung, chịu trách nhiệm hoàn toàn về Active Directory Domain Services (AD DS), cụm DNS/DHCP sẵn sàng cao, Network Policy Server (NPS) và cấu hình File Services.
        \item Thiết kế kiến trúc mạng Core, Distribution và Access mạnh mẽ, kết hợp thiết bị mạng tiêu chuẩn với định tuyến VLAN phức tạp và quy hoạch IP subnetting.
        \item Đảm bảo an ninh nghiêm ngặt trên toàn hạ tầng viễn thông bằng cách triển khai và tinh chỉnh Next-Generation Firewalls (NGFW), Web Application Firewalls (WAF), Site-to-Site VPNs bảo mật và Layer 7 Load Balancers.
        \item Thiết kế và vận hành các chiến lược Khôi phục Thảm họa (DR) và bảo vệ dữ liệu liên tục bằng Veeam Backup \& Replication, SRM, và nhân bản đa điểm (multi-site replication) nhằm tuân thủ nghiêm ngặt RPO/RTO.
        \item Tập trung hóa khả năng quan sát hệ thống thông qua việc tích hợp LibreNMS, Zabbix và Grafana để giám sát hiệu năng theo thời gian thực, tổng hợp syslog và cảnh báo SNMP cho toàn bộ máy chủ và thiết bị mạng.
        \item Thu hẹp khoảng cách giữa hạ tầng và phần mềm bằng cách duy trì các hệ thống xác thực Smart OTP (Python) và luồng phân tích dữ liệu viễn thông.
    \end{highlights}
\end{onecolentry}
\vspace{0.2 cm}

% --- ActsOne ---
\begin{twocolentry}{\textit{2024 -- Tháng 4 2025}}
    \textbf{ActsOne} \\[4pt]
    \textit{Kiến trúc sư Giải pháp / Kỹ sư Hệ thống}
\end{twocolentry}
\begin{onecolentry}
    \begin{highlights}
        \item Xây dựng nền tảng xử lý đơn hàng (Golang, gRPC, DDD) tích hợp Shopee, Lazada, TikTok với cơ chế quản lý phiên bản tồn kho và khóa phân tán (distributed locking), phục vụ đắc lực cho hoạt động thương mại số.
        \item Thiết kế và vận hành hạ tầng IDC on-premise (Proxmox VE, Kubernetes, lưu trữ NAS, mạng) với cơ chế failover HA.
        \item Triển khai chiến lược lưu trữ phân tầng hot/warm/cold trên hệ thống NAS để quản lý vòng đời dữ liệu quy mô lớn.
        \item Cấu hình nhân bản MySQL/PostgreSQL (master-slave, multi-master) để đảm bảo tính sẵn sàng cao.
        \item Kiến trúc nền tảng bảo mật zero-trust (Cloudflare ZeroTrust, đường hầm bảo mật, WAF, chống DDoS).
        \item Xây dựng hạ tầng giám sát SRE và các pipeline CI/CD (GitHub Actions, Drone CI, Kubernetes) rút ngắn time-to-market.
        \item Triển khai các cụm k8s/k3s thông qua Ansible với bộ lưu trữ bền vững Longhorn.
        \item Thiết lập mạng lưới bảo mật toàn công ty (DrayTek firewalls, Headscale VPN cùng Tailscale/WireGuard).
        \item Triển khai Authentik làm IdP trung tâm (SSO, MFA, RBAC) cho các dịch vụ nội bộ.
        \item Phát triển hệ thống API bằng Golang với tài liệu Swagger/Scalar, Redis caching và API Gateway.
    \end{highlights}
\end{onecolentry}
\vspace{0.2 cm}

% --- Kays.vn ---
\begin{twocolentry}{\textit{2023}}
    \textbf{Kays.vn} \\[4pt]
    \textit{Kỹ sư Backend}
\end{twocolentry}
\begin{onecolentry}
    \begin{highlights}
        \item Xây dựng các module backend đa nền tảng bằng Rust để xử lý giao dịch, hỗ trợ mã hóa và xử lý dữ liệu bảo mật.
        \item Thiết kế và triển khai pipeline dữ liệu thời gian thực hướng sự kiện bằng AWS EventBridge để đồng bộ giữa Supabase và hệ thống cũ, phục vụ cá nhân hóa dữ liệu khách hàng.
        \item Phát triển RESTful APIs và các dịch vụ backend cho phân hệ danh mục sản phẩm và quản lý đơn hàng.
        \item Đào tạo (mentor) 4 thực tập sinh về các mẫu kiến trúc backend và tối ưu hóa hiệu năng.
    \end{highlights}
\end{onecolentry}
\vspace{0.2 cm}

% --- IVS ---
\begin{twocolentry}{\textit{2022}}
    \textbf{IVS Technology Solutions} \\[4pt]
    \textit{Lập trình viên Backend}
\end{twocolentry}
\begin{onecolentry}
    \begin{highlights}
        \item Phát triển và bảo trì các RESTful APIs bằng SpringBoot (Java) cho các module nghiệp vụ cốt lõi.
        \item Tối ưu hóa truy vấn MongoDB/MySQL với các chiến lược lập chỉ mục và phân tích truy vấn, giúp giảm 40\% thời gian phản hồi toàn hệ thống.
        \item Thiết kế schema cơ sở dữ liệu và triển khai các tầng truy cập dữ liệu (data access layers) cho các khối lượng công việc đọc/ghi thông lượng cao.
        \item Xây dựng pipeline CI/CD bằng Jenkins để tự động hóa toàn bộ quá trình kiểm thử và triển khai.
    \end{highlights}
\end{onecolentry}

\section{Kỹ năng Công nghệ \& Sở thích}
\begin{onecolentry}
\begin{highlights}
    \item \textbf{Phân tích Dữ liệu \& MarTech:} Data Pipelines, Event-driven Architecture, Tích hợp API TMĐT (Shopee, TikTok, Lazada)
    \item \textbf{Lập trình:} Go, C/C++, Python, JavaScript, Java, Rust
    \item \textbf{Backend:} gRPC, Domain-Driven Design (DDD), Microservices, API Gateway
    \item \textbf{Hạ tầng (Infrastructure):} On-Prem IDC, Proxmox VE, Kubernetes (k8s/k3s), Telecom Infrastructure
    \item \textbf{Cơ sở dữ liệu:} MySQL/PostgreSQL (Replication, Master-Slave), MongoDB, Redis, DuckDB
    \item \textbf{Bảo mật:} Cloudflare ZeroTrust, Secure Tunnels, DrayTek Firewalls, Zero-Trust Architecture
    \item \textbf{Mạng (Networking):} VLAN, Load Balancing, Network Segmentation, VPN (Tailscale/WireGuard)
    \item \textbf{DevOps:} Docker, Ansible, GitHub Actions, Drone CI, Nginx, CI/CD Pipelines
    \item \textbf{Web:} Next.js, React.js, SpringBoot, Tauri, Electron
\end{highlights}
\end{onecolentry}

\section{Thành tựu}
\begin{onecolentry}
    \begin{highlights}
        \item Giải pháp Tốt nhất (Best Solution Award) - ActsOne (2024)
        \item Nhân viên Xuất sắc (Outstanding Employee) - Kays.vn (2023)
        \item Thực tập sinh Triển vọng (Most Promising Intern Award) - IVS Technology Solutions (2022)
        \item Học bổng Khuyến khích Học tập (Academic Excellence Scholarship) - HCMUT (2019, 2022)
        \item Giải thưởng Dự án Sinh viên (Điểm: 9.30/10) - HCMUT (2022)
    \end{highlights}
\end{onecolentry}
